\chapter{Исследование \texttt{stride} и пространственной локальности}
\label{ch:stride}

\section{Параметры исследования}

Размер массива превышает объём исследуемого кэша: $N = 10^7$, $\mathit{repeats} = 20$.
Размер строки кэша --- \texttt{128 B}.

Количество фактических обращений и нормализованная стоимость одного обращения:
$$
N_{access} = \frac{N}{\mathit{stride}} \times \mathit{repeats}, \qquad
\mathit{cycles/access} = \frac{\mathit{cycles}}{N_{access}}.
$$

L1d miss rate рассчитывается так же, как в главе~\ref{ch:working-set}.
Утилита \texttt{mperf}, как и \texttt{perf stat}, учитывает весь процесс, включая выделение и заполнение массивов.
Эта часть не зависит от \texttt{stride} и при больших \texttt{stride} сопоставима с самим циклом, поэтому из всех показателей вычитается запуск с $\mathit{repeats} = 0$.

\begin{commands}
sudo python3 scripts/stride.py
\end{commands}

\includescript{stride.py}{Измерения для разных значений \texttt{stride}}

\section{Измеренные показатели}

Показатели холостого запуска: \texttt{cycles} --- \texttt{64410983}, загрузки --- \texttt{2612952}, промахи L1d --- \texttt{204891}.

\begin{table}[H]
\caption{Зависимость показателей от \texttt{stride}}
\label{tab:stride}
\small
\begin{tabularx}{\textwidth}{|r|r|C|C|C|C|}
	\hline
	\texttt{stride} & $N_{access}$ & \texttt{cycles} & \textit{cycles/ access} & Промахи L1d & L1d miss rate, \% \\
	\hline
	\texttt{1} & \texttt{200000000} & \texttt{230063313} & \texttt{1.15} & \texttt{11096310} & \texttt{2.77} \\
	\hline
	\texttt{2} & \texttt{100000000} & \texttt{165006890} & \texttt{1.65} & \texttt{13896876} & \texttt{6.95} \\
	\hline
	\texttt{4} & \texttt{50000000} & \texttt{163639851} & \texttt{3.27} & \texttt{40323205} & \texttt{40.32} \\
	\hline
	\texttt{8} & \texttt{25000000} & \texttt{156800735} & \texttt{6.27} & \texttt{37317107} & \texttt{74.63} \\
	\hline
	\texttt{16} & \texttt{12500000} & \texttt{233745496} & \texttt{18.70} & \texttt{24352419} & \texttt{97.41} \\
	\hline
	\texttt{32} & \texttt{6250000} & \texttt{149043598} & \texttt{23.85} & \texttt{11925774} & \texttt{95.41} \\
	\hline
	\texttt{64} & \texttt{3125000} & \texttt{81764859} & \texttt{26.16} & \texttt{6106280} & \texttt{97.70} \\
	\hline
\end{tabularx}
\end{table}

\pgfplotstableread{data/stride.dat}\stridedata
\begin{figure}[H]
	\centering
	\begin{tikzpicture}
		\begin{axis}[
			width=0.95\textwidth,
			height=7cm,
			font=\small,
			xmode=log,
			log basis x=2,
			xtick={1,2,4,8,16,32,64},
			xticklabels={1,2,4,8,16,32,64},
			extra x ticks={16},
			extra x tick labels={строка кэша},
			extra x tick style={xticklabel pos=upper, grid=major, major grid style={dashed}},
			xlabel={\texttt{stride}},
			ylabel={\textit{cycles/access}},
		]
			\addplot[mark=*] table[x=stride, y=cycles_per_access]{\stridedata};
		\end{axis}
	\end{tikzpicture}
	\caption{Зависимость \textit{cycles/access} от \texttt{stride}}
	\label{fig:stride-cycles}
\end{figure}

\section{Вывод}

В пятом этапе получена зависимость cycles/access и показателей промахов кэша от stride.
До stride = 8 стоимость обращения растёт почти вдвое при каждом увеличении шага, так как из загруженной строки кэша используется всё меньше элементов.
При stride = 16 каждое обращение попадает в новую строку (128 B = 16 элементов double), промахи L1d достигают 97\% и cycles/access резко возрастает, дальше рост замедляется.
